% ======================================================================
% Working draft for IEEE ISCAS 2027
%
% IMPORTANT TEMPLATE NOTE (2026-08-26):
% The ISCAS 2027 regular-paper submission page is online, but the detailed
% 2027 author kit/template has not yet been published. This draft therefore
% follows the latest publicly available ISCAS 2026 author instructions:
% IEEE conference format, 4 pages of technical content + optional 5th page
% containing references only. Replace only the class/template wrapper if the
% 2027 author kit introduces any formatting change after submissions open.
%
% Working title intentionally avoids an acronym because "DiVE/DIVE" is already
% used by multiple video-diffusion works.
%
% Paper positioning:
%   - Algorithm novelty: Quantization-Aligned Caching (QAC). The cache indicator
%     is evaluated/calibrated in the same quantized numerical domain as the
%     recompute path, addressing decision-domain mismatch in naive Cache+Quant.
%   - Hardware novelty: a dimension-elastic row-wise PEA:
%       (1) output-axis elasticity for QK^T via 4-head x 9-lane partitioning;
%       (2) reduction-axis elasticity for PV via boundary-aware segmented reduction.
%   - Current RTL note: segmented temporal PV mode is enabled for T >= 18.
%     The Open-Sora v1.2 default T=15 short-T path must be finalized before
%     claiming default-workload PV improvement in the final paper.
% ======================================================================
% IMPLEMENTATION NOTE (2026-09-03):
% 28-nm DC synthesis values now inserted: 4.2624 mm^2, 116.372 mW,
% 1 GHz, 1296 INT8 + 36 FP16 MACs, 238.5 KB SRAM, 256 GB/s (8 DMA channels).
% Performance is modeled by a cycle-level performance simulator.
% Preliminary 7-nm normalization: 0.134 mm^2, 43.6 mW, Fmax=1.264 GHz.
% Before submission, verify whether the 28-nm area includes SRAM macros.
% ======================================================================

\documentclass[conference]{IEEEtran}

\usepackage{cite}
\usepackage{amsmath,amssymb}
\usepackage{graphicx}
\usepackage{booktabs}
\usepackage{multirow}
\usepackage{xcolor}
\usepackage{array}
\usepackage{url}

% ----------------------- Known values / placeholders --------------------
\newcommand{\ModelName}{Open-Sora v1.2}
\newcommand{\BaselineVBench}{76.77}
\newcommand{\QuantOnlyVBench}{76.79}
\newcommand{\TeaCacheVBench}{75.88}
\newcommand{\NaiveCQVBench}{76.61}
\newcommand{\QACVBench}{\textbf{[QAC-VBENCH]}}

\newcommand{\SpatialSeq}{405}
\newcommand{\TemporalSeq}{15}
\newcommand{\SpatialContexts}{30}
\newcommand{\TemporalContexts}{810}

\newcommand{\BaseSpatialQKUtil}{93.75\%}
\newcommand{\OursSpatialQKUtil}{100\%}
\newcommand{\BaseTemporalQKUtil}{41.67\%}
\newcommand{\OursTemporalQKUtil}{83.33\%}

% Algorithm placeholders: fill after QAC characterization.
\newcommand{\FPIndicatorCorr}{\textbf{[$\rho_{\mathrm{FP}}$]}}
\newcommand{\QIndicatorCorr}{\textbf{[$\rho_{\mathrm{Q}}$]}}
\newcommand{\BaseCacheRatio}{\textbf{[BASE-CACHE-RATIO]}}
\newcommand{\QACCacheRatio}{\textbf{[QAC-CACHE-RATIO]}}
\newcommand{\QACOverhead}{\textbf{[QAC-OVERHEAD]}}

% Known implementation parameters.
\newcommand{\TechNode}{TSMC 28\,nm}
\newcommand{\Frequency}{1\,GHz}
\newcommand{\SRAMCapacity}{238.5\,KB}
\newcommand{\MemoryBandwidth}{256\,GB/s}
\newcommand{\DMAChannels}{8}
\newcommand{\INTEightMACs}{1296}
\newcommand{\FPSixteenMACs}{36}
\newcommand{\AreaTwentyEight}{4.2624\,mm$^2$}
\newcommand{\PowerTwentyEight}{116.372\,mW}
\newcommand{\AreaSevenNorm}{0.134\,mm$^2$}
\newcommand{\PowerSevenNorm}{43.6\,mW}
\newcommand{\FmaxSevenNorm}{1.264\,GHz}
\newcommand{\AreaScaleFactor}{31.818}
\newcommand{\DelayScaleFactor}{1.264}
\newcommand{\PowerScaleFactor}{2.667}

% Fill after final RTL/simulator measurements.
\newcommand{\AreaOverhead}{\textbf{[AREA\%]}}
\newcommand{\PowerOverhead}{\textbf{[POWER\%]}}
\newcommand{\AttentionSpeedup}{\textbf{[ATTN-SPEEDUP$\times$]}}
\newcommand{\STDiTSpeedup}{\textbf{[STDIT-SPEEDUP$\times$]}}
\newcommand{\EtoESpeedup}{\textbf{[E2E-SPEEDUP$\times$]}}
\newcommand{\EnergyGain}{\textbf{[ENERGY-EFF.$\times$]}}
\newcommand{\QuantStepSpeedup}{\textbf{[QUANT-STEP-SPEEDUP$\times$]}}
\newcommand{\PVTemporalUtilGain}{\textbf{[PV-UTIL: XX\%$\rightarrow$YY\%]}}
\newcommand{\PVTemporalSpeedup}{\textbf{[PV-SPEEDUP$\times$]}}

% Use compact TODO markers while drafting; remove before submission.
\newcommand{\TODO}[1]{\textcolor{red}{\textbf{[TODO: #1]}}}

% Compile-safe figure placeholder. Replace with \includegraphics in final.
\newcommand{\placeholderfig}[2]{%
  \fbox{\parbox[c][#1][c]{0.96\linewidth}{\centering\small #2}}%
}

\title{Quantization-Aligned Caching and Dimension-Elastic\\
Acceleration for Video Diffusion Transformers}

% \author{
% \IEEEauthorblockN{Anonymous Author(s) / \TODO{replace with author list}}
% \IEEEauthorblockA{\TODO{affiliation and email}}
% }

\begin{document}
\maketitle

\begin{abstract}
Video Diffusion Transformers (VDiTs) achieve compelling spatiotemporal
generation quality at the cost of expensive iterative denoising. Caching and
low-bit quantization exploit inter-step and numerical redundancy,
respectively; however, their direct composition exposes two coupled
inefficiencies. First, conventional cache indicators are derived from
full-precision feature variation, whereas recomputed steps follow a quantized
trajectory, causing a \emph{decision-domain mismatch}. Second, factorized
spatial/temporal attention produces highly asymmetric matrix dimensions that
fragment fixed-width processing-element arrays. To address these problems,
this paper proposes a cross-layer co-design with
\emph{quantization-aligned caching} (QAC) and a
\emph{dimension-elastic row-wise PEA}. QAC evaluates cache relevance in the
same low-precision domain used by recomputation and calibrates the indicator
against quantized output variation, improving indicator correlation from
\FPIndicatorCorr{} to \QIndicatorCorr{} and increasing the cache ratio from
\BaseCacheRatio{} to \QACCacheRatio{} at matched generation quality. On
hardware, the default \ModelName{} workload exhibits spatial/temporal
sequence lengths of \SpatialSeq{}/\TemporalSeq{}. A four-head,
nine-lane $QK^{\mathsf T}$ mapping raises their output-lane utilization from
\BaseSpatialQKUtil{}/\BaseTemporalQKUtil{} to
\OursSpatialQKUtil{}/\OursTemporalQKUtil{}, while a boundary-aware segmented
reduction tree packs variable-length $PV$ reductions without cross-segment
interference. Evaluations on \ModelName{} achieve \QACVBench{} VBench versus
\BaselineVBench{} for the 16-bit baseline. RTL implementation in \TechNode{}
delivers \EtoESpeedup{} end-to-end speedup and \EnergyGain{} energy-efficiency
improvement with \AreaOverhead{} area overhead.
\end{abstract}

% Figure 1 should summarize the complete argument: mismatch -> method -> benefit.
\begin{figure*}[t]
  \centering
  \placeholderfig{1.42in}{
  \textbf{Fig. 1 placeholder: Two mismatches $\rightarrow$ two co-designed solutions $\rightarrow$ measured benefits.}\\[2pt]
  \textbf{Algorithmic mismatch:}
  FP-domain cache indicator $\not\equiv$ INT8 execution relevance.
  Show two adjacent steps whose FP variation and quantized-code/output
  variation disagree. $\rightarrow$ \textbf{QAC:} evaluate/calibrate the cache
  indicator in the quantized domain. $\rightarrow$ correlation
  \FPIndicatorCorr{}$\rightarrow$\QIndicatorCorr{}, cache ratio
  \BaseCacheRatio{}$\rightarrow$\QACCacheRatio{} at matched quality.\\
  \textbf{Hardware mismatch:}
  Spatial: \SpatialContexts{} contexts$\times$\SpatialSeq{} tokens;
  Temporal: \TemporalContexts{} contexts$\times$\TemporalSeq{} tokens.
  The variable token dimension maps to the output axis of $QK^{\mathsf T}$
  but the reduction axis of $PV$. $\rightarrow$ \textbf{Dimension-elastic PEA:}
  4-head$\times$9-lane $QK^{\mathsf T}$ mapping + boundary-aware segmented
  $PV$ reduction. $\rightarrow$ temporal $QK^{\mathsf T}$ utilization
  \BaseTemporalQKUtil{}$\rightarrow$\OursTemporalQKUtil{},
  E2E \EtoESpeedup{}, area +\AreaOverhead{}.
  }
  \caption{Overview of the proposed algorithm--hardware co-design. Quantization
  changes the numerical domain in which cache relevance should be evaluated,
  while factorized spatial/temporal attention changes the matrix axis on which
  parallelism becomes irregular. QAC and the dimension-elastic PEA address
  these two mismatches jointly.}
  \label{fig:overview}
\end{figure*}

\section{Introduction}
\label{sec:intro}
\IEEEPARstart{V}{ideo} Diffusion Transformers (VDiTs) have emerged as a
powerful paradigm for high-fidelity video generation
~\cite{zheng2024opensora}, but their iterative denoising incurs substantial
inference cost. Open-Sora v1.2, for example, repeatedly executes 28 spatial--temporal block pairs in a
Spatial-Temporal Diffusion Transformer (STDiT) over 30 sampling steps.
Efficient deployment therefore requires reducing both the number of STDiT
evaluations and the cost of each evaluation.

Caching and quantization exploit these two opportunities from complementary
directions. Timestep caching reuses intermediate results across adjacent
denoising steps~\cite{liu2025teacache}, while post-training quantization
replaces high-precision operators with low-bit computation
~\cite{zhao2025viditq}. Their joint use has also been explored through cache
schedule optimization, error correction, hierarchical caching, and adaptive
precision~\cite{liu2025cachequant,wu2025quantcache}. However, directly
combining timestep caching with fixed W8A8/INT8 execution leaves two
cross-layer mismatches unresolved.

First, caching and quantization operate in different \textbf{numerical
domains}. Existing cache indicators estimate the benefit of recomputation
from the variation of a full-precision feature $F_t$, whereas a quantized
model evolves according to $Q_s(F_t)$. Because quantization is nonlinear and
piecewise constant, full-precision variation need not preserve the relative
change seen by the low-precision model. An FP-domain indicator can therefore
misestimate cache relevance after quantization. Our characterization observes
only \FPIndicatorCorr{} correlation between the conventional indicator and
the actual quantized-output variation.

Second, the remaining low-precision STDiT computation exhibits a
\textbf{matrix-axis mismatch}. Under the default \ModelName{} configuration,
VAE compression and patch embedding produce $S=\SpatialSeq{}$ spatial tokens
and $T=\TemporalSeq{}$ temporal tokens. With classifier-free guidance, this
corresponds to \SpatialContexts{} spatial contexts of length \SpatialSeq{}
but \TemporalContexts{} temporal contexts of length \TemporalSeq{}. Moreover,
the varying sequence length $L$ appears on different GEMM axes,
\begin{equation}
QK^{\mathsf T}:[1,d]\!\times\![d,L],\qquad
PV:[1,L]\!\times\![L,d],
\label{eq:attn_axes}
\end{equation}
fragmenting output parallelism in $QK^{\mathsf T}$ and reduction parallelism
in $PV$. On a fixed 36$\times$36 INT8 PEA with $d=72$, spatial and temporal
$QK^{\mathsf T}$ achieve \BaseSpatialQKUtil{} and only
\BaseTemporalQKUtil{} utilization, respectively.

These observations motivate an algorithm--hardware co-design that aligns both
the cache decision and the arithmetic organization with the actual
low-precision workload. At the algorithm level, we propose
\textbf{quantization-aligned caching (QAC)}, which evaluates cache relevance
from the quantized representation and calibrates the indicator against
quantized output variation. QAC retains fixed W8A8/INT8 execution without introducing
mixed precision or an auxiliary correction network. At the architecture
level, we develop a \textbf{dimension-elastic row-wise PEA}. Its 36
$QK^{\mathsf T}$ output lanes are partitioned into four 9-lane head groups,
raising spatial/temporal lane utilization to
\OursSpatialQKUtil{}/\OursTemporalQKUtil{}, while a boundary-aware segmented
reduction tree packs variable-length $PV$ reductions across dot-product
boundaries. Stride-programmable loads provide spatial and temporal views
without explicit tensor reshaping.

On \ModelName{}, QAC improves indicator correlation from \FPIndicatorCorr{}
to \QIndicatorCorr{} and increases cache reuse from \BaseCacheRatio{} to
\QACCacheRatio{} at matched generation quality. The complete model achieves
\QACVBench{} VBench versus \BaselineVBench{} for the 16-bit baseline. RTL
evaluation in \TechNode{} shows \EtoESpeedup{} end-to-end speedup and
\EnergyGain{} higher energy efficiency with \AreaOverhead{} area overhead.

Our contributions are summarized as follows:
\begin{itemize}
  \item We identify a \textbf{decision-domain mismatch} in cache-quantized
  VDiT inference and propose \textbf{QAC}, which aligns cache decisions with
  the low-precision execution trajectory.
  \item We reveal the \textbf{matrix-axis mismatch} of factorized
  spatial/temporal attention and propose a \textbf{dimension-elastic row-wise
  PEA} that adapts output parallelism for $QK^{\mathsf T}$ and reduction
  parallelism for $PV$.
  \item We implement the co-design in RTL and evaluate it on real
  \ModelName{} tensor shapes, demonstrating \QACVBench{} VBench,
  \EtoESpeedup{} end-to-end speedup, and \EnergyGain{} energy-efficiency
  improvement.
\end{itemize}

% =========================================================================
% PAGE-BUDGET GUIDANCE:
% Page 1: Abstract + Fig. 1 + Introduction.
% Page 2: QAC + matrix-shape characterization + architecture overview.
% Page 3: dimension-elastic PEA details + experimental setup/results.
% Page 4: ablations/comparisons + conclusion.
% Page 5: references only, if needed.
% =========================================================================

\section{Quantization-Aligned Caching and Workload Characterization}
\label{sec:characterization}

\subsection{Quantization-Aligned Caching}
\label{sec:qac}

Let $F_t$ denote the low-cost feature used by a timestep-cache policy before
the expensive STDiT computation at step $t$. A conventional cache indicator
can be expressed as a normalized feature variation
\begin{equation}
d_t^{\mathrm{FP}}=
\frac{\lVert F_t-F_{t-1}\rVert_1}
     {\lVert F_{t-1}\rVert_1+\epsilon}.
\label{eq:fp_indicator}
\end{equation}
The cache predictor is calibrated to map $d_t^{\mathrm{FP}}$ to the expected
change of a full-precision model. When the actual recompute path is quantized,
this predictor is no longer aligned with the executed numerical trajectory.

QAC instead forms the indicator from the quantized representation
$\widetilde F_t=D(Q_s(F_t))$, where $Q_s$ and $D$ denote quantization and
dequantization using the calibrated scale $s$:
\begin{equation}
d_t^{\mathrm{Q}}=
\frac{\lVert \widetilde F_t-\widetilde F_{t-1}\rVert_1}
     {\lVert \widetilde F_{t-1}\rVert_1+\epsilon}.
\label{eq:q_indicator}
\end{equation}
An offline calibration run fits a lightweight mapping
$g_q(d_t^{\mathrm{Q}},t)$ to the actual output variation of the W8A8/INT8
model. At runtime, the predicted variation is accumulated using the same
thresholded reuse rule as timestep caching. Thus, QAC changes only
\emph{when} a step is recomputed; all recompute steps retain the same
hardware-friendly fixed precision. QAC uses the same calibrated quantizer as the recompute path, limiting the
additional runtime cost to \QACOverhead{}.

\begin{table}[t]
\centering
\caption{Algorithm-level quality and cache efficiency on \ModelName{}.}
\label{tab:quality}
\setlength{\tabcolsep}{3.0pt}
\begin{tabular}{lccc}
\toprule
Method & Precision & Cache ratio & VBench \\
\midrule
16-bit baseline & 16/16 & 0 & \BaselineVBench \\
ViDiT-Q~\cite{zhao2025viditq} & 8/8 & 0 & \QuantOnlyVBench \\
TeaCache~\cite{liu2025teacache} & 16/16 & [X] & \TeaCacheVBench \\
Naive Cache+Quant & 8/8 & \BaseCacheRatio & \NaiveCQVBench \\
QAC (ours) & 8/8 & \QACCacheRatio & \QACVBench \\
\bottomrule
\end{tabular}
\end{table}

\subsection{Axis-Induced Matrix Asymmetry}
\label{sec:axis_asym}

For the default workload, $T=\TemporalSeq{}$, $S=\SpatialSeq{}$, and
$d=72$. Spatial attention forms $BT=\SpatialContexts{}$ contexts of length
$S$, whereas temporal attention forms $BS=\TemporalContexts{}$ contexts of
length $T$. Hence the same attention block alternates between two reciprocal
execution regimes: temporal attention exposes $27\times$ more independent
contexts, but each context is $27\times$ shorter.

\begin{table}[t]
\centering
\caption{Default Open-Sora v1.2 attention shapes under CFG.}
\label{tab:shape}
\setlength{\tabcolsep}{4.0pt}
\begin{tabular}{lccc}
\toprule
Mode & Contexts & Seq. $L$ & Baseline $QK$ util. \\
\midrule
Spatial & \SpatialContexts & \SpatialSeq & \BaseSpatialQKUtil \\
Temporal & \TemporalContexts & \TemporalSeq & \BaseTemporalQKUtil \\
\bottomrule
\end{tabular}
\end{table}

Because our attention is row-wise fused, a complete $L\times L$ attention map
is never materialized. The hardware problem is instead \emph{shape
fragmentation}: $L$ maps to output lanes in $QK^{\mathsf T}$ but to reduction
lanes in $PV$. This observation motivates a PEA that can adapt both
parallelism axes without duplicating the arithmetic datapath.

\section{Dimension-Elastic Row-Wise Accelerator}
\label{sec:hardware}

\subsection{Architecture Overview}

\begin{figure}[t]
  \centering
  \placeholderfig{1.02in}{
  \textbf{Fig. 3 placeholder: RTL architecture.}\\
  Draw Load/AGU + IFMAP SRAM + Weight SRAM $\rightarrow$ 36$\times$36 INT8 PEA
  $\rightarrow$ row-wise softmax/VCU $\rightarrow$ OFMAP SRAM.
  Highlight two reconfiguration controls:
  (1) $QK$: 4-head$\times$9-lane partition;
  (2) $PV$: segmented reduction boundaries.
  Show stride fields selecting spatial vs temporal access.}
  \caption{The proposed accelerator reuses one row-wise PEA for $QK^{\mathsf T}$ and $PV$ while
  adapting the output and reduction dimensions independently.}
  \label{fig:arch}
\end{figure}

The baseline PEA contains 36 output lanes, each with a 36-product INT8 dot
product. Channel dimensions are naturally aligned: $d=72=2\times36$ and the
STDiT hidden dimension is $1152=32\times36$. Hence the dominant mismatch is
not the channel dimension, but the variable spatial/temporal token dimension.

\subsection{Output-Axis Elasticity for $QK^{\mathsf T}$}

A conventional mapping assigns all 36 output lanes to one attention head.
For temporal attention, only 15 lanes carry useful scores, giving
$15/36=\BaseTemporalQKUtil{}$ utilization. The proposed PEA partitions the array into
four independent 9-lane head groups:
\begin{equation}
36\ \text{lanes} \rightarrow
4\ \text{heads}\times 9\ \text{lanes/head}.
\end{equation}
For $T=15$, each head completes in two 9-lane groups, increasing temporal
utilization to $15/18=\OursTemporalQKUtil{}$. For $S=405$, $405/9=45$
exactly, eliminating the spatial tail and reaching \OursSpatialQKUtil{}.

\subsection{Reduction-Axis Elasticity for $PV$}

In $PV$, the same sequence length $L$ becomes the reduction dimension. Padding
each short temporal dot product to 36 wastes multipliers. The proposed PEA instead treats
successive dot products as a continuous product stream. Each 36-product beat
carries reduction-boundary metadata; every reduction-tree node propagates a
segment tag and only adds neighboring values with identical tags. Different
segments are compacted without being mixed, and completed reductions are
committed while an unfinished tail is accumulated into the next beat.

\begin{figure}[t]
  \centering
  \placeholderfig{0.82in}{
  \textbf{Fig. 4 placeholder: Segmented reduction.}\\
  Baseline: [15 valid | 21 pad] per dot product.\\
  Proposed co-design: continuous 36-product stream crossing boundaries:
  [tail of R0 | R1 | head of R2] $\rightarrow$ tag-aware tree
  $\rightarrow$ completed R0/R1 + partial R2.\\
  Add a small variable-$T$ latency/utilization curve after measurements.}
  \caption{Boundary-aware segmented reduction converts padding bubbles into
  useful products from adjacent dot products.}
  \label{fig:segred}
\end{figure}

% CRITICAL IMPLEMENTATION CHECK:
% Current uploaded RTL enables the segmented temporal mode only for T >= 18.
% The default Open-Sora v1.2 T=15 path must be completed (e.g., a short-T
% two-reduction packing mode) before replacing the placeholders below.
For the final RTL, report both multiplier utilization and cycle reduction over
a sweep of temporal lengths:
\begin{equation}
T\in\{\textbf{[5, 10, 15, 20, 25, 30, 36]}\}.
\end{equation}
The default $T=15$ point should report \PVTemporalUtilGain{} and
\PVTemporalSpeedup{} only after the short-$T$ mode is functionally verified.

\subsection{Stride-Programmable Spatial/Temporal Access}

STDiT alternates logical $[T,S,C]$ and $[S,T,C]$ traversal. The proposed design does not
physically reshape the activation tensor. Instead, the load instructions
program the AXI/SRAM address stride to select the required logical axis. This
mechanism is presented as a supporting feature rather than a standalone
contribution; its purpose is to feed both PEA modes without extra data
movement.

\section{Evaluation}
\label{sec:eval}

\subsection{Experimental Setup}

\textbf{Model and quality.}
We evaluate \ModelName{} using its default 240p, 9:16, 51-frame,
30-step configuration and report VBench~\cite{huang2024vbench}.
\TODO{state prompt count, random seeds, cache threshold, calibration set,
and whether VBench is averaged over all official prompts or a subset.}

\textbf{Hardware implementation and performance modeling.}
The proposed accelerator is synthesized with Synopsys Design Compiler (DC)
using a \TechNode{} standard-cell library and a target clock frequency of
\Frequency{}. DC is used to obtain the synthesized area, power, and timing
results, including the maximum achievable clock frequency. The accelerator
integrates a $36\times36$ INT8 PEA (\INTEightMACs{} INT8 MACs),
\FPSixteenMACs{} FP16 MACs in the VCU, and \SRAMCapacity{} of on-chip SRAM.
Eight DMA channels collectively transfer 256\,B/cycle, corresponding to
\MemoryBandwidth{} at 1\,GHz.

Full-workload performance is evaluated using our cycle-level performance
simulator. The simulator models the execution cycles of the compute engines,
on-chip data movement, DMA transfers, and scheduling behavior of the proposed
architecture. Reported latency and throughput are derived from the simulated
cycle counts at the evaluated clock frequency, whereas area, power, and
frequency are taken from DC synthesis.
\TODO{add the exact DC version, supply voltage, and power-activity assumption
if required in the final manuscript.}

\subsection{Implementation and Technology Normalization}
\label{subsec:hw_eval}

Table~\ref{tab:hw_comparison} summarizes the DC synthesis results and the
technology-normalized comparison with ViDA~\cite{ding2025vida}. At 28\,nm,
DC reports an area of \AreaTwentyEight{} and a power of
\PowerTwentyEight{} at 1\,GHz. Following the scaling methodology adopted by
ViDA, we normalize the 28-nm synthesis results to 7\,nm using area, delay,
and power scaling factors of \AreaScaleFactor{}, \DelayScaleFactor{}, and
\PowerScaleFactor{}, respectively.
This gives a normalized area of \AreaSevenNorm{}, power of
\PowerSevenNorm{}, and an estimated $F_{\max}$ of \FmaxSevenNorm{}.

For cross-accelerator performance comparison, we nevertheless use 1\,GHz and
256\,GB/s, matching ViDA's evaluation setting rather than exploiting the
higher normalized $F_{\max}$. ViDA reports 640 FP16 MACs and 2.03\,mm$^2$ at
7\,nm, but does not report accelerator power. Because the two designs use
different MAC counts and precisions, Table~\ref{tab:hw_comparison} is used to
state the implementation context; speedup claims are reported separately
under explicitly matched frequency and bandwidth assumptions.

% IMPORTANT BEFORE FINAL SUBMISSION:
% Confirm whether 4.2624 mm^2 is logic/core-only area or already includes
% SRAM macros. If SRAM is not included, do not present 0.134 mm^2 as a total
% chip area. Following ViDA's methodology, scale logic separately and add a
% target-node SRAM macro area (e.g., CACTI/compiler characterization).
\begin{table}[t]
\centering
\caption{Implementation parameters and 7-nm normalization.}
\label{tab:hw_comparison}
\footnotesize
\setlength{\tabcolsep}{2.4pt}
\renewcommand{\arraystretch}{1.05}
\resizebox{\columnwidth}{!}{%
\begin{tabular}{@{}lccc@{}}
\toprule
\textbf{Metric} & \textbf{Ours--28 nm} & \textbf{Ours--7 nm Norm.} & \textbf{ViDA~\cite{ding2025vida}} \\
\midrule
Frequency & 1.0\,GHz & 1.0\,GHz$^{a}$ & 1.0\,GHz \\
MACs & 1296 INT8 + 36 FP16 & 1296 INT8 + 36 FP16 & 640 FP16 \\
SRAM & 238.5\,KB & 238.5\,KB & N/R \\
Memory BW & 256\,GB/s & 256\,GB/s & 256\,GB/s \\
Area & 4.2624\,mm$^2$ & 0.134\,mm$^2$$^{b}$ & 2.03\,mm$^2$ \\
Power & 116.372\,mW & 43.6\,mW & N/R \\
\bottomrule
\end{tabular}%
}
\vspace{1pt}
\parbox{\columnwidth}{\scriptsize
$^{a}$Technology-normalized $F_{\max}$ is 1.264\,GHz; 1\,GHz is used for
comparison. $^{b}$Preliminary scaling of the reported 28-nm area; SRAM must be
normalized separately if it is not included in the synthesis area.}
\end{table}

\subsection{Microarchitectural Efficiency}

\begin{table}[t]
\centering
\caption{Ablation of dimension-elastic scheduling.}
\label{tab:hardware}
\setlength{\tabcolsep}{3.0pt}
\begin{tabular}{lcccc}
\toprule
Design & Area & Power & Attn. Lat. & E2E \\
\midrule
Baseline PEA & [X] & [X] & 1.00$\times$ & 1.00$\times$ \\
+ QK elasticity & [X] & [X] & [X] & [X] \\
+ PV elasticity & [X] & [X] & [X] & [X] \\
Full design & [X] & [X] & \AttentionSpeedup & \EtoESpeedup \\
\bottomrule
\end{tabular}
\end{table}

The final result should emphasize measured latency rather than utilization
alone. \TODO{report, at $T=15$, the temporal QK latency reduction, PV latency
reduction, combined attention speedup, and incremental area/power overhead.}

\subsection{End-to-End Co-Design Ablation}

\begin{table}[t]
\centering
\caption{End-to-end contribution breakdown.}
\label{tab:ablation}
\setlength{\tabcolsep}{3.0pt}
\begin{tabular}{lccc}
\toprule
Configuration & VBench & Latency & Speedup \\
\midrule
16-bit dense & \BaselineVBench & [X] & 1.00$\times$ \\
Quant only & \QuantOnlyVBench & [X] & [X] \\
Naive Cache+Quant & \NaiveCQVBench & [X] & [X] \\
+ QAC & \QACVBench & [X] & [X] \\
+ dimension-elastic PEA & \QACVBench & [X] & \EtoESpeedup \\
\bottomrule
\end{tabular}
\end{table}

This ablation separates the gain from numerical-domain alignment (QAC) from
the gain of the dimension-elastic PEA. The comparison therefore attributes
end-to-end speedup to the reduction in recomputed denoising steps and to the
per-step hardware acceleration independently.

\section{Conclusion}
\label{sec:conclusion}

This paper presents an algorithm--hardware co-design for efficient video
diffusion inference. At the algorithm level, QAC addresses the
decision-domain mismatch introduced by naively combining timestep caching
with low-precision execution, aligning cache decisions with the quantized
trajectory. At the hardware level, a dimension-elastic row-wise PEA addresses
the matrix-axis asymmetry of spatial/temporal attention through output-lane
partitioning for $QK^{\mathsf T}$ and boundary-aware segmented reduction for
$PV$. Evaluated on \ModelName{}, the proposed design achieves
\QACVBench{} VBench versus \BaselineVBench{} for the 16-bit baseline and
\EtoESpeedup{} end-to-end speedup with \AreaOverhead{} area overhead. These
results demonstrate that numerical-domain alignment and dimension-aware
microarchitecture are complementary for efficient low-precision video
diffusion inference.

\bibliographystyle{IEEEtran}
\bibliography{references}

\end{document}

% ======================================================================
% Working draft for IEEE ISCAS 2027
%
% IMPORTANT TEMPLATE NOTE (2026-08-26):
% The ISCAS 2027 regular-paper submission page is online, but the detailed
% 2027 author kit/template has not yet been published. This draft therefore
% follows the latest publicly available ISCAS 2026 author instructions:
% IEEE conference format, 4 pages of technical content + optional 5th page
% containing references only. Replace only the class/template wrapper if the
% 2027 author kit introduces any formatting change after submissions open.
%
% Working title intentionally avoids an acronym because "DiVE/DIVE" is already
% used by multiple video-diffusion works.
%
% Paper positioning:
%   - Algorithm novelty: Quantization-Aligned Caching (QAC). The cache indicator
%     is evaluated/calibrated in the same quantized numerical domain as the
%     recompute path, addressing decision-domain mismatch in naive Cache+Quant.
%   - Hardware novelty: a dimension-elastic row-wise PEA:
%       (1) output-axis elasticity for QK^T via 4-head x 9-lane partitioning;
%       (2) reduction-axis elasticity for PV via boundary-aware segmented reduction.
%   - Current RTL note: segmented temporal PV mode is enabled for T >= 18.
%     The Open-Sora v1.2 default T=15 short-T path must be finalized before
%     claiming default-workload PV improvement in the final paper.
% ======================================================================
% IMPLEMENTATION NOTE (2026-09-03):
% 28-nm measured/synthesized values now inserted: 4.2624 mm^2, 116.372 mW,
% 1 GHz, 1296 INT8 + 36 FP16 MACs, 238.5 KB SRAM, 256 GB/s (8 DMA channels).
% Preliminary 7-nm normalization: 0.134 mm^2, 43.6 mW, Fmax=1.264 GHz.
% Before submission, verify whether the 28-nm area includes SRAM macros.
% ======================================================================

\documentclass[conference]{IEEEtran}

\usepackage{cite}
\usepackage{amsmath,amssymb}
\usepackage{graphicx}
\usepackage{booktabs}
\usepackage{multirow}
\usepackage{xcolor}
\usepackage{array}
\usepackage{url}

% ----------------------- Known values / placeholders --------------------
\newcommand{\ModelName}{Open-Sora v1.2}
\newcommand{\BaselineVBench}{76.77}
\newcommand{\QuantOnlyVBench}{76.79}
\newcommand{\TeaCacheVBench}{75.88}
\newcommand{\NaiveCQVBench}{76.61}
\newcommand{\QACVBench}{\textbf{[QAC-VBENCH]}}

\newcommand{\SpatialSeq}{405}
\newcommand{\TemporalSeq}{15}
\newcommand{\SpatialContexts}{30}
\newcommand{\TemporalContexts}{810}

\newcommand{\BaseSpatialQKUtil}{93.75\%}
\newcommand{\OursSpatialQKUtil}{100\%}
\newcommand{\BaseTemporalQKUtil}{41.67\%}
\newcommand{\OursTemporalQKUtil}{83.33\%}

% Algorithm placeholders: fill after QAC characterization.
\newcommand{\FPIndicatorCorr}{\textbf{[$\rho_{\mathrm{FP}}$]}}
\newcommand{\QIndicatorCorr}{\textbf{[$\rho_{\mathrm{Q}}$]}}
\newcommand{\BaseCacheRatio}{\textbf{[BASE-CACHE-RATIO]}}
\newcommand{\QACCacheRatio}{\textbf{[QAC-CACHE-RATIO]}}
\newcommand{\QACOverhead}{\textbf{[QAC-OVERHEAD]}}

% Known implementation parameters.
\newcommand{\TechNode}{TSMC 28\,nm}
\newcommand{\Frequency}{1\,GHz}
\newcommand{\SRAMCapacity}{238.5\,KB}
\newcommand{\MemoryBandwidth}{256\,GB/s}
\newcommand{\DMAChannels}{8}
\newcommand{\INTEightMACs}{1296}
\newcommand{\FPSixteenMACs}{36}
\newcommand{\AreaTwentyEight}{4.2624\,mm$^2$}
\newcommand{\PowerTwentyEight}{116.372\,mW}
\newcommand{\AreaSevenNorm}{0.134\,mm$^2$}
\newcommand{\PowerSevenNorm}{43.6\,mW}
\newcommand{\FmaxSevenNorm}{1.264\,GHz}
\newcommand{\AreaScaleFactor}{31.818}
\newcommand{\DelayScaleFactor}{1.264}
\newcommand{\PowerScaleFactor}{2.667}

% Fill after final RTL/simulator measurements.
\newcommand{\AreaOverhead}{\textbf{[AREA\%]}}
\newcommand{\PowerOverhead}{\textbf{[POWER\%]}}
\newcommand{\AttentionSpeedup}{\textbf{[ATTN-SPEEDUP$\times$]}}
\newcommand{\STDiTSpeedup}{\textbf{[STDIT-SPEEDUP$\times$]}}
\newcommand{\EtoESpeedup}{\textbf{[E2E-SPEEDUP$\times$]}}
\newcommand{\EnergyGain}{\textbf{[ENERGY-EFF.$\times$]}}
\newcommand{\QuantStepSpeedup}{\textbf{[QUANT-STEP-SPEEDUP$\times$]}}
\newcommand{\PVTemporalUtilGain}{\textbf{[PV-UTIL: XX\%$\rightarrow$YY\%]}}
\newcommand{\PVTemporalSpeedup}{\textbf{[PV-SPEEDUP$\times$]}}

% Use compact TODO markers while drafting; remove before submission.
\newcommand{\TODO}[1]{\textcolor{red}{\textbf{[TODO: #1]}}}

% Compile-safe figure placeholder. Replace with \includegraphics in final.
\newcommand{\placeholderfig}[2]{%
  \fbox{\parbox[c][#1][c]{0.96\linewidth}{\centering\small #2}}%
}

\title{Quantization-Aligned Caching and Dimension-Elastic\\
Acceleration for Video Diffusion Transformers}

% \author{
% \IEEEauthorblockN{Anonymous Author(s) / \TODO{replace with author list}}
% \IEEEauthorblockA{\TODO{affiliation and email}}
% }

\begin{document}
\maketitle

\begin{abstract}
Video Diffusion Transformers (VDiTs) achieve compelling spatiotemporal
generation quality at the cost of expensive iterative denoising. Caching and
low-bit quantization exploit inter-step and numerical redundancy,
respectively; however, their direct composition exposes two coupled
inefficiencies. First, conventional cache indicators are derived from
full-precision feature variation, whereas recomputed steps follow a quantized
trajectory, causing a \emph{decision-domain mismatch}. Second, factorized
spatial/temporal attention produces highly asymmetric matrix dimensions that
fragment fixed-width processing-element arrays. To address these problems,
this paper proposes a cross-layer co-design with
\emph{quantization-aligned caching} (QAC) and a
\emph{dimension-elastic row-wise PEA}. QAC evaluates cache relevance in the
same low-precision domain used by recomputation and calibrates the indicator
against quantized output variation, improving indicator correlation from
\FPIndicatorCorr{} to \QIndicatorCorr{} and increasing the cache ratio from
\BaseCacheRatio{} to \QACCacheRatio{} at matched generation quality. On
hardware, the default \ModelName{} workload exhibits spatial/temporal
sequence lengths of \SpatialSeq{}/\TemporalSeq{}. A four-head,
nine-lane $QK^{\mathsf T}$ mapping raises their output-lane utilization from
\BaseSpatialQKUtil{}/\BaseTemporalQKUtil{} to
\OursSpatialQKUtil{}/\OursTemporalQKUtil{}, while a boundary-aware segmented
reduction tree packs variable-length $PV$ reductions without cross-segment
interference. Evaluations on \ModelName{} achieve \QACVBench{} VBench versus
\BaselineVBench{} for the 16-bit baseline. RTL implementation in \TechNode{}
delivers \EtoESpeedup{} end-to-end speedup and \EnergyGain{} energy-efficiency
improvement with \AreaOverhead{} area overhead.
\end{abstract}

% Figure 1 should summarize the complete argument: mismatch -> method -> benefit.
\begin{figure*}[t]
  \centering
  \placeholderfig{1.42in}{
  \textbf{Fig. 1 placeholder: Two mismatches $\rightarrow$ two co-designed solutions $\rightarrow$ measured benefits.}\\[2pt]
  \textbf{Algorithmic mismatch:}
  FP-domain cache indicator $\not\equiv$ INT8 execution relevance.
  Show two adjacent steps whose FP variation and quantized-code/output
  variation disagree. $\rightarrow$ \textbf{QAC:} evaluate/calibrate the cache
  indicator in the quantized domain. $\rightarrow$ correlation
  \FPIndicatorCorr{}$\rightarrow$\QIndicatorCorr{}, cache ratio
  \BaseCacheRatio{}$\rightarrow$\QACCacheRatio{} at matched quality.\\
  \textbf{Hardware mismatch:}
  Spatial: \SpatialContexts{} contexts$\times$\SpatialSeq{} tokens;
  Temporal: \TemporalContexts{} contexts$\times$\TemporalSeq{} tokens.
  The variable token dimension maps to the output axis of $QK^{\mathsf T}$
  but the reduction axis of $PV$. $\rightarrow$ \textbf{Dimension-elastic PEA:}
  4-head$\times$9-lane $QK^{\mathsf T}$ mapping + boundary-aware segmented
  $PV$ reduction. $\rightarrow$ temporal $QK^{\mathsf T}$ utilization
  \BaseTemporalQKUtil{}$\rightarrow$\OursTemporalQKUtil{},
  E2E \EtoESpeedup{}, area +\AreaOverhead{}.
  }
  \caption{Overview of the proposed algorithm--hardware co-design. Quantization
  changes the numerical domain in which cache relevance should be evaluated,
  while factorized spatial/temporal attention changes the matrix axis on which
  parallelism becomes irregular. QAC and the dimension-elastic PEA address
  these two mismatches jointly.}
  \label{fig:overview}
\end{figure*}

\section{Introduction}
\label{sec:intro}
\IEEEPARstart{V}{ideo} Diffusion Transformers (VDiTs) have emerged as a
powerful paradigm for high-fidelity video generation
~\cite{zheng2024opensora}, but their iterative denoising incurs substantial
inference cost. Open-Sora v1.2, for example, repeatedly executes 28 spatial--temporal block pairs in a
Spatial-Temporal Diffusion Transformer (STDiT) over 30 sampling steps.
Efficient deployment therefore requires reducing both the number of STDiT
evaluations and the cost of each evaluation.

Caching and quantization exploit these two opportunities from complementary
directions. Timestep caching reuses intermediate results across adjacent
denoising steps~\cite{liu2025teacache}, while post-training quantization
replaces high-precision operators with low-bit computation
~\cite{zhao2025viditq}. Their joint use has also been explored through cache
schedule optimization, error correction, hierarchical caching, and adaptive
precision~\cite{liu2025cachequant,wu2025quantcache}. However, directly
combining timestep caching with fixed W8A8/INT8 execution leaves two
cross-layer mismatches unresolved.

First, caching and quantization operate in different \textbf{numerical
domains}. Existing cache indicators estimate the benefit of recomputation
from the variation of a full-precision feature $F_t$, whereas a quantized
model evolves according to $Q_s(F_t)$. Because quantization is nonlinear and
piecewise constant, full-precision variation need not preserve the relative
change seen by the low-precision model. An FP-domain indicator can therefore
misestimate cache relevance after quantization. Our characterization observes
only \FPIndicatorCorr{} correlation between the conventional indicator and
the actual quantized-output variation.

Second, the remaining low-precision STDiT computation exhibits a
\textbf{matrix-axis mismatch}. Under the default \ModelName{} configuration,
VAE compression and patch embedding produce $S=\SpatialSeq{}$ spatial tokens
and $T=\TemporalSeq{}$ temporal tokens. With classifier-free guidance, this
corresponds to \SpatialContexts{} spatial contexts of length \SpatialSeq{}
but \TemporalContexts{} temporal contexts of length \TemporalSeq{}. Moreover,
the varying sequence length $L$ appears on different GEMM axes,
\begin{equation}
QK^{\mathsf T}:[1,d]\!\times\![d,L],\qquad
PV:[1,L]\!\times\![L,d],
\label{eq:attn_axes}
\end{equation}
fragmenting output parallelism in $QK^{\mathsf T}$ and reduction parallelism
in $PV$. On a fixed 36$\times$36 INT8 PEA with $d=72$, spatial and temporal
$QK^{\mathsf T}$ achieve \BaseSpatialQKUtil{} and only
\BaseTemporalQKUtil{} utilization, respectively.

These observations motivate an algorithm--hardware co-design that aligns both
the cache decision and the arithmetic organization with the actual
low-precision workload. At the algorithm level, we propose
\textbf{quantization-aligned caching (QAC)}, which evaluates cache relevance
from the quantized representation and calibrates the indicator against
quantized output variation. QAC retains fixed W8A8/INT8 execution without introducing
mixed precision or an auxiliary correction network. At the architecture
level, we develop a \textbf{dimension-elastic row-wise PEA}. Its 36
$QK^{\mathsf T}$ output lanes are partitioned into four 9-lane head groups,
raising spatial/temporal lane utilization to
\OursSpatialQKUtil{}/\OursTemporalQKUtil{}, while a boundary-aware segmented
reduction tree packs variable-length $PV$ reductions across dot-product
boundaries. Stride-programmable loads provide spatial and temporal views
without explicit tensor reshaping.

On \ModelName{}, QAC improves indicator correlation from \FPIndicatorCorr{}
to \QIndicatorCorr{} and increases cache reuse from \BaseCacheRatio{} to
\QACCacheRatio{} at matched generation quality. The complete model achieves
\QACVBench{} VBench versus \BaselineVBench{} for the 16-bit baseline. RTL
evaluation in \TechNode{} shows \EtoESpeedup{} end-to-end speedup and
\EnergyGain{} higher energy efficiency with \AreaOverhead{} area overhead.

Our contributions are summarized as follows:
\begin{itemize}
  \item We identify a \textbf{decision-domain mismatch} in cache-quantized
  VDiT inference and propose \textbf{QAC}, which aligns cache decisions with
  the low-precision execution trajectory.
  \item We reveal the \textbf{matrix-axis mismatch} of factorized
  spatial/temporal attention and propose a \textbf{dimension-elastic row-wise
  PEA} that adapts output parallelism for $QK^{\mathsf T}$ and reduction
  parallelism for $PV$.
  \item We implement the co-design in RTL and evaluate it on real
  \ModelName{} tensor shapes, demonstrating \QACVBench{} VBench,
  \EtoESpeedup{} end-to-end speedup, and \EnergyGain{} energy-efficiency
  improvement.
\end{itemize}

% =========================================================================
% PAGE-BUDGET GUIDANCE:
% Page 1: Abstract + Fig. 1 + Introduction.
% Page 2: QAC + matrix-shape characterization + architecture overview.
% Page 3: dimension-elastic PEA details + experimental setup/results.
% Page 4: ablations/comparisons + conclusion.
% Page 5: references only, if needed.
% =========================================================================

\section{Quantization-Aligned Caching and Workload Characterization}
\label{sec:characterization}

\subsection{Quantization-Aligned Caching}
\label{sec:qac}

Let $F_t$ denote the low-cost feature used by a timestep-cache policy before
the expensive STDiT computation at step $t$. A conventional cache indicator
can be expressed as a normalized feature variation
\begin{equation}
d_t^{\mathrm{FP}}=
\frac{\lVert F_t-F_{t-1}\rVert_1}
     {\lVert F_{t-1}\rVert_1+\epsilon}.
\label{eq:fp_indicator}
\end{equation}
The cache predictor is calibrated to map $d_t^{\mathrm{FP}}$ to the expected
change of a full-precision model. When the actual recompute path is quantized,
this predictor is no longer aligned with the executed numerical trajectory.

QAC instead forms the indicator from the quantized representation
$\widetilde F_t=D(Q_s(F_t))$, where $Q_s$ and $D$ denote quantization and
dequantization using the calibrated scale $s$:
\begin{equation}
d_t^{\mathrm{Q}}=
\frac{\lVert \widetilde F_t-\widetilde F_{t-1}\rVert_1}
     {\lVert \widetilde F_{t-1}\rVert_1+\epsilon}.
\label{eq:q_indicator}
\end{equation}
An offline calibration run fits a lightweight mapping
$g_q(d_t^{\mathrm{Q}},t)$ to the actual output variation of the W8A8/INT8
model. At runtime, the predicted variation is accumulated using the same
thresholded reuse rule as timestep caching. Thus, QAC changes only
\emph{when} a step is recomputed; all recompute steps retain the same
hardware-friendly fixed precision. QAC uses the same calibrated quantizer as the recompute path, limiting the
additional runtime cost to \QACOverhead{}.

\begin{table}[t]
\centering
\caption{Algorithm-level quality and cache efficiency on \ModelName{}.}
\label{tab:quality}
\setlength{\tabcolsep}{3.0pt}
\begin{tabular}{lccc}
\toprule
Method & Precision & Cache ratio & VBench \\
\midrule
16-bit baseline & 16/16 & 0 & \BaselineVBench \\
ViDiT-Q~\cite{zhao2025viditq} & 8/8 & 0 & \QuantOnlyVBench \\
TeaCache~\cite{liu2025teacache} & 16/16 & [X] & \TeaCacheVBench \\
Naive Cache+Quant & 8/8 & \BaseCacheRatio & \NaiveCQVBench \\
QAC (ours) & 8/8 & \QACCacheRatio & \QACVBench \\
\bottomrule
\end{tabular}
\end{table}

\subsection{Axis-Induced Matrix Asymmetry}
\label{sec:axis_asym}

For the default workload, $T=\TemporalSeq{}$, $S=\SpatialSeq{}$, and
$d=72$. Spatial attention forms $BT=\SpatialContexts{}$ contexts of length
$S$, whereas temporal attention forms $BS=\TemporalContexts{}$ contexts of
length $T$. Hence the same attention block alternates between two reciprocal
execution regimes: temporal attention exposes $27\times$ more independent
contexts, but each context is $27\times$ shorter.

\begin{table}[t]
\centering
\caption{Default Open-Sora v1.2 attention shapes under CFG.}
\label{tab:shape}
\setlength{\tabcolsep}{4.0pt}
\begin{tabular}{lccc}
\toprule
Mode & Contexts & Seq. $L$ & Baseline $QK$ util. \\
\midrule
Spatial & \SpatialContexts & \SpatialSeq & \BaseSpatialQKUtil \\
Temporal & \TemporalContexts & \TemporalSeq & \BaseTemporalQKUtil \\
\bottomrule
\end{tabular}
\end{table}

Because our attention is row-wise fused, a complete $L\times L$ attention map
is never materialized. The hardware problem is instead \emph{shape
fragmentation}: $L$ maps to output lanes in $QK^{\mathsf T}$ but to reduction
lanes in $PV$. This observation motivates a PEA that can adapt both
parallelism axes without duplicating the arithmetic datapath.

\section{Dimension-Elastic Row-Wise Accelerator}
\label{sec:hardware}

\subsection{Architecture Overview}

\begin{figure}[t]
  \centering
  \placeholderfig{1.02in}{
  \textbf{Fig. 3 placeholder: RTL architecture.}\\
  Draw Load/AGU + IFMAP SRAM + Weight SRAM $\rightarrow$ 36$\times$36 INT8 PEA
  $\rightarrow$ row-wise softmax/VCU $\rightarrow$ OFMAP SRAM.
  Highlight two reconfiguration controls:
  (1) $QK$: 4-head$\times$9-lane partition;
  (2) $PV$: segmented reduction boundaries.
  Show stride fields selecting spatial vs temporal access.}
  \caption{The proposed accelerator reuses one row-wise PEA for $QK^{\mathsf T}$ and $PV$ while
  adapting the output and reduction dimensions independently.}
  \label{fig:arch}
\end{figure}

The baseline PEA contains 36 output lanes, each with a 36-product INT8 dot
product. Channel dimensions are naturally aligned: $d=72=2\times36$ and the
STDiT hidden dimension is $1152=32\times36$. Hence the dominant mismatch is
not the channel dimension, but the variable spatial/temporal token dimension.

\subsection{Output-Axis Elasticity for $QK^{\mathsf T}$}

A conventional mapping assigns all 36 output lanes to one attention head.
For temporal attention, only 15 lanes carry useful scores, giving
$15/36=\BaseTemporalQKUtil{}$ utilization. The proposed PEA partitions the array into
four independent 9-lane head groups:
\begin{equation}
36\ \text{lanes} \rightarrow
4\ \text{heads}\times 9\ \text{lanes/head}.
\end{equation}
For $T=15$, each head completes in two 9-lane groups, increasing temporal
utilization to $15/18=\OursTemporalQKUtil{}$. For $S=405$, $405/9=45$
exactly, eliminating the spatial tail and reaching \OursSpatialQKUtil{}.

\subsection{Reduction-Axis Elasticity for $PV$}

In $PV$, the same sequence length $L$ becomes the reduction dimension. Padding
each short temporal dot product to 36 wastes multipliers. The proposed PEA instead treats
successive dot products as a continuous product stream. Each 36-product beat
carries reduction-boundary metadata; every reduction-tree node propagates a
segment tag and only adds neighboring values with identical tags. Different
segments are compacted without being mixed, and completed reductions are
committed while an unfinished tail is accumulated into the next beat.

\begin{figure}[t]
  \centering
  \placeholderfig{0.82in}{
  \textbf{Fig. 4 placeholder: Segmented reduction.}\\
  Baseline: [15 valid | 21 pad] per dot product.\\
  Proposed co-design: continuous 36-product stream crossing boundaries:
  [tail of R0 | R1 | head of R2] $\rightarrow$ tag-aware tree
  $\rightarrow$ completed R0/R1 + partial R2.\\
  Add a small variable-$T$ latency/utilization curve after measurements.}
  \caption{Boundary-aware segmented reduction converts padding bubbles into
  useful products from adjacent dot products.}
  \label{fig:segred}
\end{figure}

% CRITICAL IMPLEMENTATION CHECK:
% Current uploaded RTL enables the segmented temporal mode only for T >= 18.
% The default Open-Sora v1.2 T=15 path must be completed (e.g., a short-T
% two-reduction packing mode) before replacing the placeholders below.
For the final RTL, report both multiplier utilization and cycle reduction over
a sweep of temporal lengths:
\begin{equation}
T\in\{\textbf{[5, 10, 15, 20, 25, 30, 36]}\}.
\end{equation}
The default $T=15$ point should report \PVTemporalUtilGain{} and
\PVTemporalSpeedup{} only after the short-$T$ mode is functionally verified.

\subsection{Stride-Programmable Spatial/Temporal Access}

STDiT alternates logical $[T,S,C]$ and $[S,T,C]$ traversal. The proposed design does not
physically reshape the activation tensor. Instead, the load instructions
program the AXI/SRAM address stride to select the required logical axis. This
mechanism is presented as a supporting feature rather than a standalone
contribution; its purpose is to feed both PEA modes without extra data
movement.

\section{Evaluation}
\label{sec:eval}

\subsection{Experimental Setup}

\textbf{Model and quality.}
We evaluate \ModelName{} using its default 240p, 9:16, 51-frame,
30-step configuration and report VBench~\cite{huang2024vbench}.
\TODO{state prompt count, random seeds, cache threshold, calibration set,
and whether VBench is averaged over all official prompts or a subset.}

\textbf{Hardware implementation.}
The proposed accelerator is synthesized in \TechNode{} at \Frequency{}.
It integrates a $36\times36$ INT8 PEA (\INTEightMACs{} INT8 MACs),
\FPSixteenMACs{} FP16 MACs in the VCU, and \SRAMCapacity{} of on-chip SRAM.
Eight DMA channels collectively transfer 256\,B/cycle, corresponding to
\MemoryBandwidth{} at 1\,GHz. \TODO{add synthesis tool/version, voltage,
power methodology, and whether the reported area/power are post-synthesis or
post-layout.}

\subsection{Implementation and Technology Normalization}
\label{subsec:hw_eval}

Table~\ref{tab:hw_comparison} summarizes the physical implementation and the
technology-normalized comparison with ViDA~\cite{ding2025vida}. The 28-nm
implementation occupies \AreaTwentyEight{} and consumes \PowerTwentyEight{}
at 1\,GHz. Following the scaling methodology adopted by ViDA, we normalize
the 28-nm results to 7\,nm using area, delay, and power scaling factors of
\AreaScaleFactor{}, \DelayScaleFactor{}, and \PowerScaleFactor{}, respectively.
This gives a normalized area of \AreaSevenNorm{}, power of
\PowerSevenNorm{}, and an estimated $F_{\max}$ of \FmaxSevenNorm{}.

For cross-accelerator performance comparison, we nevertheless use 1\,GHz and
256\,GB/s, matching ViDA's evaluation setting rather than exploiting the
higher normalized $F_{\max}$. ViDA reports 640 FP16 MACs and 2.03\,mm$^2$ at
7\,nm, but does not report accelerator power. Because the two designs use
different MAC counts and precisions, Table~\ref{tab:hw_comparison} is used to
state the implementation context; speedup claims are reported separately
under explicitly matched frequency and bandwidth assumptions.

% IMPORTANT BEFORE FINAL SUBMISSION:
% Confirm whether 4.2624 mm^2 is logic/core-only area or already includes
% SRAM macros. If SRAM is not included, do not present 0.134 mm^2 as a total
% chip area. Following ViDA's methodology, scale logic separately and add a
% target-node SRAM macro area (e.g., CACTI/compiler characterization).
\begin{table}[t]
\centering
\caption{Implementation parameters and 7-nm normalization.}
\label{tab:hw_comparison}
\footnotesize
\setlength{\tabcolsep}{2.4pt}
\renewcommand{\arraystretch}{1.05}
\resizebox{\columnwidth}{!}{%
\begin{tabular}{@{}lccc@{}}
\toprule
\textbf{Metric} & \textbf{Ours--28 nm} & \textbf{Ours--7 nm Norm.} & \textbf{ViDA~\cite{ding2025vida}} \\
\midrule
Frequency & 1.0\,GHz & 1.0\,GHz$^{a}$ & 1.0\,GHz \\
MACs & 1296 INT8 + 36 FP16 & 1296 INT8 + 36 FP16 & 640 FP16 \\
SRAM & 238.5\,KB & 238.5\,KB & N/R \\
Memory BW & 256\,GB/s & 256\,GB/s & 256\,GB/s \\
Area & 4.2624\,mm$^2$ & 0.134\,mm$^2$$^{b}$ & 2.03\,mm$^2$ \\
Power & 116.372\,mW & 43.6\,mW & N/R \\
\bottomrule
\end{tabular}%
}
\vspace{1pt}
\parbox{\columnwidth}{\scriptsize
$^{a}$Technology-normalized $F_{\max}$ is 1.264\,GHz; 1\,GHz is used for
comparison. $^{b}$Preliminary scaling of the reported 28-nm area; SRAM must be
normalized separately if it is not included in the synthesis area.}
\end{table}

\subsection{Microarchitectural Efficiency}

\begin{table}[t]
\centering
\caption{Ablation of dimension-elastic scheduling.}
\label{tab:hardware}
\setlength{\tabcolsep}{3.0pt}
\begin{tabular}{lcccc}
\toprule
Design & Area & Power & Attn. Lat. & E2E \\
\midrule
Baseline PEA & [X] & [X] & 1.00$\times$ & 1.00$\times$ \\
+ QK elasticity & [X] & [X] & [X] & [X] \\
+ PV elasticity & [X] & [X] & [X] & [X] \\
Full design & [X] & [X] & \AttentionSpeedup & \EtoESpeedup \\
\bottomrule
\end{tabular}
\end{table}

The final result should emphasize measured latency rather than utilization
alone. \TODO{report, at $T=15$, the temporal QK latency reduction, PV latency
reduction, combined attention speedup, and incremental area/power overhead.}

\subsection{End-to-End Co-Design Ablation}

\begin{table}[t]
\centering
\caption{End-to-end contribution breakdown.}
\label{tab:ablation}
\setlength{\tabcolsep}{3.0pt}
\begin{tabular}{lccc}
\toprule
Configuration & VBench & Latency & Speedup \\
\midrule
16-bit dense & \BaselineVBench & [X] & 1.00$\times$ \\
Quant only & \QuantOnlyVBench & [X] & [X] \\
Naive Cache+Quant & \NaiveCQVBench & [X] & [X] \\
+ QAC & \QACVBench & [X] & [X] \\
+ dimension-elastic PEA & \QACVBench & [X] & \EtoESpeedup \\
\bottomrule
\end{tabular}
\end{table}

This ablation separates the gain from numerical-domain alignment (QAC) from
the gain of the dimension-elastic PEA. The comparison therefore attributes
end-to-end speedup to the reduction in recomputed denoising steps and to the
per-step hardware acceleration independently.

\section{Conclusion}
\label{sec:conclusion}

This paper presents an algorithm--hardware co-design for efficient video
diffusion inference. At the algorithm level, QAC addresses the
decision-domain mismatch introduced by naively combining timestep caching
with low-precision execution, aligning cache decisions with the quantized
trajectory. At the hardware level, a dimension-elastic row-wise PEA addresses
the matrix-axis asymmetry of spatial/temporal attention through output-lane
partitioning for $QK^{\mathsf T}$ and boundary-aware segmented reduction for
$PV$. Evaluated on \ModelName{}, the proposed design achieves
\QACVBench{} VBench versus \BaselineVBench{} for the 16-bit baseline and
\EtoESpeedup{} end-to-end speedup with \AreaOverhead{} area overhead. These
results demonstrate that numerical-domain alignment and dimension-aware
microarchitecture are complementary for efficient low-precision video
diffusion inference.

\bibliographystyle{IEEEtran}
\bibliography{references}

\end{document}

% This document is compiled using pdfLaTeX
% You can switch XeLaTeX/pdfLaTeX/LaTeX/LuaLaTeX in Settings

\documentclass{article}

\title{Title}
\author{Your Name}
\date{\today}

\begin{document}

\maketitle

\section{First Section}

Hello World!

\end{document}
